\documentclass{article}
\usepackage{graphicx} % Required for inserting images
\documentclass{article}
\usepackage{graphicx} % Required for inserting images
\usepackage{amsmath}
\usepackage{amssymb}
\usepackage{algorithm}
\usepackage{algpseudocode}

\title{Sutton and Barton Exercises Chapter 13}
\author{danieltocco0 }
\date{April 2026}

\begin{document}

\maketitle

\section{Chapter 13}

\textbf{Exercise 13.1.} Use your knowledge of the gridworld and its dynamics to determine an
\textit{exact} symbolic expression for the optimal probability of selecting the \textbf{right} action in
Example 13.1.

\bigskip

\textbf{Solution.}

Let $p \in [0,1]$ denote the probability of selecting \textbf{right} in all three nonterminal states
(since all states appear identical under the feature representation). Let $\gamma = 1$ and
reward $r = -1$ per step.

\bigskip

\textbf{State layout:} $S_1$ (start) $\to$ $S_2$ $\to$ $S_3$ $\to$ $G$ (terminal).

\begin{itemize}
  \item In $S_1$: \textbf{right} moves right, \textbf{left} causes no movement.
  \item In $S_2$: actions are \emph{reversed} — \textbf{right} moves left, \textbf{left} moves right.
  \item In $S_3$: \textbf{right} moves right (to $G$), \textbf{left} moves left (to $S_2$).
\end{itemize}

\bigskip

\textbf{Expected steps to goal.}

From $S_3$: with probability $p$ we reach $G$, with probability $1-p$ we go to $S_2$.

From $S_2$: with probability $p$ we go to $S_1$, with probability $1-p$ we go to $S_3$.

From $S_1$: with probability $p$ we move to $S_2$, with probability $1-p$ we stay in $S_1$.

\bigskip

Let $v_i$ denote the (negative) value, i.e., expected number of steps to $G$, from state $S_i$.
Since reward is $-1$ per step, $J(\theta) = v_{\pi_\theta}(S_1) = -v_1$.

\bigskip

Setting up the system: since \textbf{left} from $S_1$ causes no movement, the expected steps
from $S_1$ satisfy:

\begin{equation}
  v_1 = 1 + (1-p)\,v_1 + p\,v_2
\end{equation}

\begin{equation}
  v_2 = 1 + p\,v_1 + (1-p)\,v_3
\end{equation}

\begin{equation}
  v_3 = 1 + (1-p)\,v_2
\end{equation}

From equation (1):
\[
  p\,v_1 = 1 + p\,v_2 \implies v_1 = \frac{1}{p} + v_2.
\]

Substituting into (2):
\[
  v_2 = 1 + p\!\left(\frac{1}{p} + v_2\right) + (1-p)\,v_3
       = 2 + p\,v_2 + (1-p)\,v_3,
\]
\[
  (1-p)\,v_2 = 2 + (1-p)\,v_3 \implies v_2 = \frac{2}{1-p} + v_3.
\]

Substituting $v_3$ from (3):
\[
  v_3 = 1 + (1-p)\!\left(\frac{2}{1-p} + v_3\right) = 3 + (1-p)\,v_3,
\]
\[
  p\,v_3 = 3 \implies v_3 = \frac{3}{p}.
\]

Back-substituting:
\[
  v_2 = \frac{2}{1-p} + \frac{3}{p}, \qquad
  v_1 = \frac{1}{p} + \frac{2}{1-p} + \frac{3}{p} = \frac{4}{p} + \frac{2}{1-p}.
\]

The performance objective is:
\[
  J(\theta) = -v_1 = -\frac{4}{p} - \frac{2}{1-p}.
\]

\bigskip

\textbf{Optimizing over $p$.} Differentiating and setting $\frac{dJ}{dp} = 0$:
\[
  \frac{dJ}{dp} = \frac{4}{p^2} - \frac{2}{(1-p)^2} = 0,
\]
\[
  \frac{4}{p^2} = \frac{2}{(1-p)^2} \implies \frac{(1-p)^2}{p^2} = \frac{1}{2},
\]
\[
  \frac{1-p}{p} = \frac{1}{\sqrt{2}} \implies 1 - p = \frac{p}{\sqrt{2}},
\]
\[
  1 = p\!\left(1 + \frac{1}{\sqrt{2}}\right) = p \cdot \frac{\sqrt{2}+1}{\sqrt{2}}.
\]

\[
  \boxed{p^* = \frac{\sqrt{2}}{\sqrt{2}+1} = \sqrt{2}\,(\sqrt{2}-1) = 2 - \sqrt{2} \approx 0.586.}
\]

This matches the optimal value shown in Example 13.1 (approximately $0.59$), with $J(p^*) \approx -11.6$.
\\\\
\textbf{Exercise 13.2.} Generalize the box on page 199 (the on-policy distribution in episodic
tasks), the policy gradient theorem (13.5), the proof of the policy gradient theorem (page 325),
and the steps leading to the REINFORCE update equation (13.8), so that (13.8) ends up with
a factor of $\gamma^t$ and thus aligns with the general algorithm given in the pseudocode.

\bigskip

\textbf{Solution.}

\bigskip

\textbf{Step 1: On-policy distribution with discounting.}

Without discounting, the expected number of visits to state $s$ in a single episode satisfies
(Equation 9.2):
\[
  \eta(s) = h(s) + \sum_{\bar{s}} \eta(\bar{s}) \sum_a \pi(a|\bar{s})\, p(s|\bar{s}, a),
\]
where $h(s)$ is the probability that an episode begins in $s$. With discounting ($\gamma < 1$),
future visitations are worth less, so we include a factor of $\gamma$ in the second term:
\[
  \eta(s) = h(s) + \gamma \sum_{\bar{s}} \eta(\bar{s}) \sum_a \pi(a|\bar{s})\, p(s|\bar{s}, a).
\]
The discounted on-policy distribution is then:
\[
  \mu(s) = \frac{\eta(s)}{\sum_{s'} \eta(s')}, \qquad \text{for all } s \in \mathcal{S},
\]
where now $\eta(s)$ encodes discounted visitation counts, so states visited later in the episode
receive geometrically smaller weight.

\bigskip

\textbf{Step 2: Policy gradient theorem with discounting.}

The unrolling in the proof of the policy gradient theorem produces transition probabilities
$\Pr(s \to x, k, \pi)$ — the probability of reaching $x$ from $s$ in exactly $k$ steps under $\pi$.
With discounting, each $k$-step contribution is weighted by $\gamma^k$:
\[
  \nabla v_\pi(s) = \sum_{x \in \mathcal{S}} \sum_{k=0}^{\infty} \gamma^k \Pr(s \to x, k, \pi)
  \sum_a \nabla \pi(a|x, \boldsymbol{\theta})\, q_\pi(x, a).
\]
Since $\nabla J(\boldsymbol{\theta}) = \nabla v_\pi(s_0)$, and unrolling from $s_0$ means $k$-step
transitions correspond to timestep $t = k$, we obtain the discounted policy gradient theorem:
\[
  \nabla J(\boldsymbol{\theta}) \propto \sum_s \mu(s) \sum_a q_\pi(s, a)\, \nabla \pi(a|s, \boldsymbol{\theta}),
  \tag{13.5$'$}
\]
where $\mu(s)$ is now the \emph{discounted} on-policy distribution defined in Step 1. The
structure of the theorem is unchanged; only the meaning of $\mu(s)$ has been updated to
incorporate $\gamma^t$ weighting.

\bigskip

\textbf{Step 3: Proof of the discounted policy gradient theorem.}

Following the same steps as the undiscounted proof, but carrying $\gamma^k$ through the
unrolling:
\begin{align*}
  \nabla J(\boldsymbol{\theta})
  &= \nabla v_\pi(s_0) \\
  &= \sum_s \left(\sum_{k=0}^{\infty} \gamma^k \Pr(s_0 \to s, k, \pi)\right)
     \sum_a \nabla\pi(a|s, \boldsymbol{\theta})\, q_\pi(s, a) \\
  &= \sum_s \eta(s) \sum_a \nabla\pi(a|s, \boldsymbol{\theta})\, q_\pi(s, a) \\
  &= \sum_{s'} \eta(s') \sum_s \frac{\eta(s)}{\sum_{s'}\eta(s')}
     \sum_a \nabla\pi(a|s, \boldsymbol{\theta})\, q_\pi(s, a) \\
  &= \sum_{s'} \eta(s') \sum_s \mu(s)
     \sum_a \nabla\pi(a|s, \boldsymbol{\theta})\, q_\pi(s, a) \\
  &\propto \sum_s \mu(s) \sum_a \nabla\pi(a|s, \boldsymbol{\theta})\, q_\pi(s, a). \qquad \text{(Q.E.D.)}
\end{align*}
The key difference from the undiscounted case is that $\eta(s) = \sum_{k=0}^\infty \gamma^k
\Pr(s_0 \to s, k, \pi)$ now discounts visits by when they occur in the episode.

\bigskip

\textbf{Step 4: REINFORCE update with $\gamma^t$.}

Starting from (13.5$'$) and writing the gradient as an expectation:
\[
  \nabla J(\boldsymbol{\theta})
  \propto \mathbb{E}_\pi\!\left[
    \sum_a q_\pi(S_t, a)\, \nabla\pi(a|S_t, \boldsymbol{\theta})
  \right].
\]
Introducing the importance sampling ratio as before:
\[
  = \mathbb{E}_\pi\!\left[
    \sum_a \pi(a|S_t, \boldsymbol{\theta})\, q_\pi(S_t, a)
    \frac{\nabla\pi(a|S_t, \boldsymbol{\theta})}{\pi(a|S_t, \boldsymbol{\theta})}
  \right].
\]
Replacing the sum over $a$ by a sample $A_t \sim \pi$:
\[
  = \mathbb{E}_\pi\!\left[
    q_\pi(S_t, A_t) \frac{\nabla\pi(A_t|S_t, \boldsymbol{\theta})}{\pi(A_t|S_t, \boldsymbol{\theta})}
  \right].
\]
Replacing $q_\pi(S_t, A_t)$ by the sample return $G_t$ (since $\mathbb{E}_\pi[G_t | S_t, A_t]
= q_\pi(S_t, A_t)$), and incorporating the $\gamma^t$ factor from the discounted $\eta(s)$
— which weights timestep $t$ by $\gamma^t$ relative to $t=0$ — we obtain:
\[
  \nabla J(\boldsymbol{\theta})
  \propto \mathbb{E}_\pi\!\left[
    \gamma^t G_t \frac{\nabla\pi(A_t|S_t, \boldsymbol{\theta})}{\pi(A_t|S_t, \boldsymbol{\theta})}
  \right].
\]
This yields the discounted REINFORCE update:
\[
  \boxed{\boldsymbol{\theta}_{t+1} = \boldsymbol{\theta}_t
  + \alpha\, \gamma^t G_t \frac{\nabla\pi(A_t|S_t, \boldsymbol{\theta}_t)}{\pi(A_t|S_t, \boldsymbol{\theta}_t)}.}
  \tag{13.8$'$}
\]

The $\gamma^t$ factor ensures that parameter updates earlier in the episode receive greater
weight, consistent with the discounted performance objective $J(\boldsymbol{\theta}) = v_\pi(s_0)$
under $\gamma < 1$. As $\gamma \to 1$, the factor $\gamma^t \to 1$ for all $t$ and we recover
the undiscounted update (13.8).
\\\\
\textbf{Exercise 13.3.} In Section 13.1 we considered policy parameterizations using the
soft-max in action preferences (13.2) with linear action preferences (13.3). For this
parameterization, prove that the eligibility vector is
\[
  \nabla \ln \pi(a|s, \boldsymbol{\theta}) = \mathbf{x}(s, a) - \sum_b \pi(b|s, \boldsymbol{\theta})\,\mathbf{x}(s, b),
  \tag{13.9}
\]
using the definitions and elementary calculus.

\bigskip

\textbf{Solution.}

The softmax policy with linear action preferences is:
\[
  \pi(a|s, \boldsymbol{\theta}) = \frac{e^{h(s,a,\boldsymbol{\theta})}}{\sum_b e^{h(s,b,\boldsymbol{\theta})}},
  \qquad h(s, a, \boldsymbol{\theta}) = \boldsymbol{\theta}^\top \mathbf{x}(s, a).
\]

Rather than applying the quotient rule directly, we first take the logarithm, which
separates the numerator and denominator:
\[
  \ln \pi(a|s, \boldsymbol{\theta})
  = h(s, a, \boldsymbol{\theta}) - \ln \sum_b e^{h(s,b,\boldsymbol{\theta})}
  = \boldsymbol{\theta}^\top \mathbf{x}(s,a) - \ln \sum_b e^{\boldsymbol{\theta}^\top \mathbf{x}(s,b)}.
\]

Now we differentiate with respect to $\boldsymbol{\theta}$. The first term is immediate:
\[
  \nabla\left[\boldsymbol{\theta}^\top \mathbf{x}(s,a)\right] = \mathbf{x}(s,a).
\]

For the second term, we apply the chain rule. The derivative of $\ln f(\boldsymbol{\theta})$
is $\nabla f(\boldsymbol{\theta}) / f(\boldsymbol{\theta})$, so:
\[
  \nabla \ln \sum_b e^{\boldsymbol{\theta}^\top \mathbf{x}(s,b)}
  = \frac{\nabla \sum_b e^{\boldsymbol{\theta}^\top \mathbf{x}(s,b)}}{\sum_b e^{\boldsymbol{\theta}^\top \mathbf{x}(s,b)}}
  = \frac{\sum_b e^{\boldsymbol{\theta}^\top \mathbf{x}(s,b)}\, \mathbf{x}(s,b)}{\sum_b e^{\boldsymbol{\theta}^\top \mathbf{x}(s,b)}}.
\]

Recognizing the softmax in the final expression:
\[
  = \sum_b \frac{e^{\boldsymbol{\theta}^\top \mathbf{x}(s,b)}}{\sum_{b'} e^{\boldsymbol{\theta}^\top \mathbf{x}(s,b')}}\,\mathbf{x}(s,b)
  = \sum_b \pi(b|s, \boldsymbol{\theta})\,\mathbf{x}(s,b).
\]

Combining both terms:
\[
  \boxed{\nabla \ln \pi(a|s, \boldsymbol{\theta}) = \mathbf{x}(s,a) - \sum_b \pi(b|s,\boldsymbol{\theta})\,\mathbf{x}(s,b).}
\]

The second term is the expected feature vector under the current policy $\pi(\cdot|s,\boldsymbol{\theta})$.
The eligibility vector therefore measures how much the feature of the chosen action $a$
deviates from the policy-weighted average feature across all actions — actions whose
features differ most from the average receive the largest gradient signal.
\\\\
\title{\textbf{Exercise 13.4 --- Gaussian Policy Eligibility Vectors}}
Show that for the Gaussian policy parameterization, the eligibility vector
$\nabla \ln \pi(a \mid s, \boldsymbol{\theta})$ has the following two parts:

\[
\nabla \ln \pi(a \mid s, \boldsymbol{\theta}_\mu)
= \frac{1}{\sigma(s,\boldsymbol{\theta})^2}\bigl(a - \mu(s,\boldsymbol{\theta})\bigr)\,\mathbf{x}_\mu(s)
\]

\[
\nabla \ln \pi(a \mid s, \boldsymbol{\theta}_\sigma)
= \left(\frac{(a - \mu(s,\boldsymbol{\theta}))^2}{\sigma(s,\boldsymbol{\theta})^2} - 1\right)\mathbf{x}_\sigma(s)
\]

where $\mu(s,\boldsymbol{\theta}) = \boldsymbol{\theta}_\mu^\top \mathbf{x}_\mu(s)$ and
$\sigma(s,\boldsymbol{\theta}) = e^{\boldsymbol{\theta}_\sigma^\top \mathbf{x}_\sigma(s)}$.

\subsection*{Solution}

Taking the log of the Gaussian policy:
\[
\ln \pi(a \mid s, \boldsymbol{\theta})
= -\ln\sigma - \ln\sqrt{2\pi} - \frac{(a-\mu)^2}{2\sigma^2}
\]

\subsubsection*{Part 1: Gradient with respect to $\boldsymbol{\theta}_\mu$}

\[
\frac{\partial \ln\pi}{\partial \mu}
= \frac{a - \mu}{\sigma^2}
\]

Since $\mu = \boldsymbol{\theta}_\mu^\top \mathbf{x}_\mu(s)$, we have $\nabla_{\boldsymbol{\theta}_\mu}\mu = \mathbf{x}_\mu(s)$, so by the chain rule:

\[
\nabla_{\boldsymbol{\theta}_\mu} \ln\pi
= \frac{a - \mu(s,\boldsymbol{\theta})}{\sigma(s,\boldsymbol{\theta})^2}\,\mathbf{x}_\mu(s)
\]

\subsubsection*{Part 2: Gradient with respect to $\boldsymbol{\theta}_\sigma$}

\[
\frac{\partial \ln\pi}{\partial \sigma}
= -\frac{1}{\sigma} + \frac{(a-\mu)^2}{\sigma^3}
= \frac{1}{\sigma}\left(\frac{(a-\mu)^2}{\sigma^2} - 1\right)
\]

Since $\sigma = e^{\boldsymbol{\theta}_\sigma^\top \mathbf{x}_\sigma(s)}$, we have $\nabla_{\boldsymbol{\theta}_\sigma}\sigma = \sigma\,\mathbf{x}_\sigma(s)$, so by the chain rule:

\[
\nabla_{\boldsymbol{\theta}_\sigma} \ln\pi
= \frac{1}{\sigma}\left(\frac{(a-\mu)^2}{\sigma^2} - 1\right) \cdot \sigma\,\mathbf{x}_\sigma(s)
= \left(\frac{(a - \mu(s,\boldsymbol{\theta}))^2}{\sigma(s,\boldsymbol{\theta})^2} - 1\right)\mathbf{x}_\sigma(s)
\]
\subsection*{Exercise 13.5 --- Bernoulli-Logistic Unit}

A \textit{Bernoulli-logistic unit} is a stochastic neuron-like unit with feature vector
$\mathbf{x}(s)$, binary action $A_t \in \{0, 1\}$, and action preferences given by
$h(s, 1, \boldsymbol{\theta}) - h(s, 0, \boldsymbol{\theta}) = \boldsymbol{\theta}^\top \mathbf{x}(s)$.

\subsubsection*{Part (a): Show $P_t = \pi(1 \mid S_t, \boldsymbol{\theta}_t) = \sigma(\boldsymbol{\theta}^\top \mathbf{x}(s))$}

Applying the softmax distribution (Eq.\ 13.2) to action $a = 1$:
\[
\pi(1 \mid s, \boldsymbol{\theta})
= \frac{e^{h(s,1,\boldsymbol{\theta})}}{e^{h(s,0,\boldsymbol{\theta})} + e^{h(s,1,\boldsymbol{\theta})}}
\]

Dividing numerator and denominator by $e^{h(s,1,\boldsymbol{\theta})}$:
\[
= \frac{1}{e^{h(s,0,\boldsymbol{\theta}) - h(s,1,\boldsymbol{\theta})} + 1}
= \frac{1}{1 + e^{-\boldsymbol{\theta}^\top \mathbf{x}(s)}}
= \sigma\!\left(\boldsymbol{\theta}^\top \mathbf{x}(s)\right) \qquad \square
\]

\subsubsection*{Part (b): Monte-Carlo REINFORCE Update}

The general REINFORCE update is:
\[
\boldsymbol{\theta}_{t+1} = \boldsymbol{\theta}_t + \alpha G_t \nabla \ln \pi(A_t \mid S_t, \boldsymbol{\theta}_t)
\]

Using the eligibility derived in part (c):
\[
\boldsymbol{\theta}_{t+1} = \boldsymbol{\theta}_t + \alpha G_t (A_t - P_t)\,\mathbf{x}(S_t)
\]

\subsubsection*{Part (c): Eligibility $\nabla \ln \pi(a \mid s, \boldsymbol{\theta})$}

Let $P = \pi(1 \mid s, \boldsymbol{\theta}) = \sigma(\boldsymbol{\theta}^\top \mathbf{x}(s))$,
so $\pi(0 \mid s, \boldsymbol{\theta}) = 1 - P$.

\textbf{Case $a = 1$:}
\[
\nabla \ln \pi(1 \mid s, \boldsymbol{\theta})
= \frac{\nabla P}{P}
\]
Using the logistic derivative $\nabla P = P(1-P)\,\mathbf{x}(s)$:
\[
= \frac{P(1-P)\,\mathbf{x}(s)}{P} = (1 - P)\,\mathbf{x}(s)
\]

\textbf{Case $a = 0$:}
\[
\nabla \ln \pi(0 \mid s, \boldsymbol{\theta})
= \frac{-\nabla P}{1 - P}
= \frac{-P(1-P)\,\mathbf{x}(s)}{1-P} = -P\,\mathbf{x}(s)
\]

Both cases unify as $(a - P)\,\mathbf{x}(s)$, verified by substitution:
\begin{align*}
a = 1: &\quad (1 - P)\,\mathbf{x}(s) \checkmark \\
a = 0: &\quad (0 - P)\,\mathbf{x}(s) = -P\,\mathbf{x}(s) \checkmark
\end{align*}

Therefore:
\[
\boxed{\nabla \ln \pi(a \mid s, \boldsymbol{\theta}) = \bigl(a - \pi(1 \mid s, \boldsymbol{\theta})\bigr)\,\mathbf{x}(s)}
\]

\hfill$\square$

\end{document}
